%% =====================================================================
%%  RUTA DE CIENCIA  —  Post #01
%%  Pilar: MITO vs. EVIDENCIA  (Lunes, Semana 1)
%%  Tema:  "Solo usas el 10% de tu cerebro"
%%
%%  Compilar:  pdflatex post-01-cerebro-10-porciento.tex
%%  Exportar:  pdftoppm -scale-to-x 1080 -png post-01-cerebro-10-porciento.pdf out/post01
%% =====================================================================
\documentclass[10pt]{beamer}
\usepackage{ruta-de-ciencia}

% ---- Metadatos de este post ----------------------------------------
\renewcommand{\rdchandle}{@rutadeciencia}
\renewcommand{\rdccategory}{MITO vs.\ EVIDENCIA}
\renewcommand{\rdctotalslides}{8}

\begin{document}

% --- 1 · PORTADA ----------------------------------------------------
\rdccover{MITO}
  {``Solo usas el\\ 10\% de tu\\ cerebro''}
  {Lo repiten las películas, los gurús del ``potencial oculto'' y ese
   post motivacional. La neurociencia dice otra cosa.}

% --- 2 · EL MITO ----------------------------------------------------
\rdccontent{2}{\faFilm}{EL MITO}
  {Lo has oído mil veces}
  {Películas como \emph{Lucy} o \emph{Sin límites} lo dan por hecho.
   Coaches lo usan para venderte que ``desbloquees'' tu potencial.
   Suena inspirador\dots\ y es falso.}

% --- 3 · LA PREGUNTA ------------------------------------------------
\rdccontent{3}{\faQuestion}{LA PREGUNTA INCÓMODA}
  {¿De dónde salió ese 10\%?}
  {De ningún estudio. Nadie midió jamás un cerebro ``90\% apagado''.
   Es una cifra que se repite sola, sin fuente que la sostenga.
   Primera señal de alarma.}

% --- 4 · EVIDENCIA 1 ------------------------------------------------
\rdccontent{4}{\faBrain}{LA EVIDENCIA · 1}
  {Te vemos el cerebro entero}
  {Las técnicas de neuroimagen (fMRI, PET) muestran actividad en
   prácticamente todo el cerebro a lo largo del día. No existe un 90\%
   en silencio esperando a ser ``activado''.}

% --- 5 · EVIDENCIA 2 ------------------------------------------------
\rdccontent{5}{\faBolt}{LA EVIDENCIA · 2}
  {La evolución no derrocha}
  {El cerebro es el 2\% de tu peso pero gasta \mbox{$\sim$20\%} de tu
   energía. Mantener un 90\% inútil sería un lujo imposible. Y una lesión
   en casi cualquier zona tiene consecuencias: no hay partes ``de sobra''.}

% --- 6 · EL ORIGEN --------------------------------------------------
\rdccontent{6}{\faHistory}{EL ORIGEN DEL MITO}
  {¿Y por qué lo creemos?}
  {Nació de malentendidos: una frase de William James sobre nuestro
   ``potencial'', lecturas erróneas sobre las células gliales y la
   autoayuda del siglo XX. Se volvió leyenda por lo bonito que suena.}

% --- 7 · LO QUE SÍ ES CIERTO ---------------------------------------
\rdccontent{7}{\faLightbulb}{ENTONCES\dots}
  {Usas el 100\%}
  {Solo que no todo a la vez. Igual que no tensas todos los músculos
   al mismo tiempo, tu cerebro activa distintas regiones según la tarea.
   Es eficiente, no perezoso.}

% --- 8 · CIERRE / VEREDICTO ----------------------------------------
\rdcclosing{8}
  {Usas todo tu cerebro. El ``10\%'' es puro mito.}
  {La próxima vez que te prometan ``desbloquear el 90\% restante'',
   ya sabes qué responder: pídeles la fuente.}
  {Beyerstein (1999), \emph{Mind Myths}\,; Society for Neuroscience\,;
   revisiones de neuroimagen fMRI/PET.}

\end{document}
